\chapter*{Preface}

This monograph is concerned with a single definition and its consequences.
The definition (\cref{def:cat}) concerns tuples consisting of an initial object,
a terminal object, and a family of objects of a fixed category, together with
the morphisms that join them. We call such a tuple a \emph{cat}, and the
subject developed from it \emph{cat theory}.

The reader is entitled to know at the outset how much is at stake. The answer
is: very little. The head of a cat is an initial object, and the tail is a
terminal object. Initial and terminal objects are unique up to unique
isomorphism, and any two parallel morphisms out of an initial object, or into
a terminal object, coincide. It follows that the head and the tail of a cat
carry no information beyond their existence, that the morphisms joining them
to the legs are uniquely determined, and that every commutative diagram we
impose on these morphisms commutes for trivial reasons. Cat theory is
therefore the theory of small families of objects of a category, a subject
which has been understood for a long time. We prove this reduction in
\cref{thm:fundamental} and use it freely afterwards.

Why, then, write a monograph? The purpose of the present work is
architectural rather than substantive. We have attempted to construct a
complete, internally consistent theory, with its own terminology, its own
named conditions, and its own correspondence theorem, while adhering
throughout to the standards of rigour one would expect of a research
monograph in category theory. Every result is stated with its hypotheses,
every proof is given in full, every type signature has been checked, and every
proposed generalization that fails is accompanied by a counterexample. Where
a feline statement is a special case or a reformulation of a classical one, we
say so explicitly. We do not claim novelty for any result in this book.

The reader who wishes to evaluate the theory quickly may consult the
dictionary in \cref{app:dictionary}, which records for each feline notion the
classical notion to which it is equivalent.

\subsection*{Organization}

\Cref{ch:intro} gives an overview and states the principal results.
\Cref{ch:prelim} collects the categorical background. \Cref{ch:cats}
introduces $\cC$-cats, their morphisms, and the category $\Fel(\cC)$, and
proves the Fundamental Theorem identifying $\Fel(\cC)$ with $\Fam(\cC)$.
\Cref{ch:anatomy} treats kittens, litters, cats satisfying the Chloe
Condition, and cats lacking a head, a tail or legs.
\Cref{ch:functors} studies feline functors and natural transformations.
\Cref{ch:limits,ch:adjunctions,ch:yoneda} concern limits and colimits,
adjunctions and equivalences, and Yoneda-type results respectively; the last
of these introduces Margaret Compatibility.
\Cref{ch:monoidal} treats monoidal and enriched cats, \cref{ch:fibrations}
feline fibrations and several derived constructions, and \cref{ch:higher}
iterated, gaited and two-dimensional generalizations.
\Cref{ch:representation} proves representation and decomposition theorems and
introduces Gideon Decompositions. \Cref{ch:cmg} proves the CMG Correspondence
Theorem relating the Chloe Condition, Margaret Compatibility and Gideon
Decompositions, together with a recognition theorem for categories of cats.
\Cref{ch:conclusion} concludes.

\subsection*{Logical dependence}

\Cref{ch:prelim,ch:cats} are required for everything that follows.
\Cref{ch:anatomy,ch:functors,ch:limits,ch:adjunctions} may be read in order;
later chapters depend on them. \Cref{ch:cmg} uses
\cref{ch:anatomy,ch:limits,ch:adjunctions,ch:yoneda,ch:representation} and, in its
last section, \cref{ch:monoidal}.
